% Chapter 5 of The Hydrogen Atoms of Deep Learning: signal propagation in wide random networks.
\newcommand{\signalpropV}{\mathcal{V}}
\newcommand{\signalpropC}{\mathcal{C}}
\newcommand{\signalpropG}{\widehat{G}}

\casechapter{Order and Chaos at Initialization}{Signal propagation in wide random networks}{Infinitely wide random networks}{ch:signalprop}
\begin{refsection}[chapters/05-signalprop.bib]

\begin{chaptersummary}
A deep network trains only if its random initialization is scaled with care, and the deepest trainable network is set by a depth scale that diverges on a curve in the plane of weight and bias variances. Infinitely wide random networks explain this. We prove that their pre-activations are Gaussian processes whose covariance obeys a layer-to-layer recursion, and that the recursion has two phases: an ordered phase, where all inputs become perfectly correlated at a geometric rate $\chi_1$ per layer, and a chaotic phase, where nearby inputs decorrelate; the depth scale is $-1/\log$ of the slope of the correlation map at its stable fixed point, and it diverges like $1/\abs{\chi_1-1}$. At He initialization ReLU networks sit exactly at the boundary and decorrelate only polynomially, $1-c^\ell\sim 9\pi^2/(2\ell^2)$. The model is mean-field theory: it is exact when width is infinite at fixed depth, it has no training and no architecture, and at finite width its fluctuations grow with depth over width, so that criticality on average says little about the typical network once $L$ is comparable to $m$.
\end{chaptersummary}

\section{The phenomenon: a narrow window of trainable initializations}\label{sec:signalprop:phenomenon}

Before a network learns anything, its weights are random, and signals must already pass through it: forward, so that the output depends on the input, and backward, so that every layer receives a gradient. Each layer multiplies the signal by a random matrix and bends it through a nonlinearity. Composing $L$ such maps either shrinks or amplifies, and generically it does so exponentially in $L$.

For gradients through time in recurrent networks the problem was identified early. \Textcite{Hochreiter1991}, in a diploma thesis written in German, analyzed how the error signal propagated backward through a recurrent network decays or blows up exponentially with the number of time steps, and \textcite{BengioEtAl1994} showed that a recurrent network that stores information robustly must have vanishing gradients over long time lags. These are the \term{vanishing and exploding gradients}\index{vanishing and exploding gradients}.\sidenote{An unrolled recurrent network is a deep feedforward network whose layers share their weights, so the two problems are one problem, with depth playing the role of time.}

For depth the practice came first. \Textcite{GlorotBengio2010} found that deep sigmoid networks with the then-standard initialization drove their top hidden layer into saturation early in training, and proposed to draw the weights of a layer with $m_{\mathrm{in}}$ inputs and $m_{\mathrm{out}}$ outputs with variance $2/(m_{\mathrm{in}}+m_{\mathrm{out}})$, a compromise between keeping the variance of the activations and that of the back-propagated gradients constant in a linear network. \Textcite{HeEtAl2015} redid the calculation for rectifiers: a ReLU passes on half of the second moment of its input, so the weight variance must be $2/m_{\mathrm{in}}$ to keep the variance of the responses constant. With this initialization they trained a $30$-layer rectifier network from scratch, while with the variance of \textcite{GlorotBengio2010} the same network stalled.

\Textcite{PooleEtAl2016} turned the variance-counting arguments into a theory. Following two inputs through a deep random network of large width, they found two phases. In the \term{ordered phase}\index{ordered phase} every pair of inputs is mapped to nearly the same representation; in the \term{chaotic phase}\index{chaotic phase} nearby inputs are driven apart, and a smooth curve of inputs is mapped to a curve whose length and curvature grow exponentially with depth. \Textcite{SchoenholzEtAl2017} read off depth scales from the same theory: the number of layers over which the correlation between two inputs relaxes, and over which gradient norms grow or decay. The correlation depth scale diverges on the boundary between the phases, the \term{edge of chaos}\index{edge of chaos}, and their experiments put the limit of trainability at a fixed multiple of it.

\begin{observation}[\textcite{SchoenholzEtAl2017}]\label[observation]{obs:signalprop:depth}
Fully connected $\tanh$ networks with Gaussian weights of variance $\sigma_w^2/m$ and biases of variance $\sigma_b^2$, were trained by stochastic gradient methods on MNIST and on CIFAR-10 for many values of the depth $L$ and of $(\sigma_w^2,\sigma_b^2)$. Networks trained when $L\lesssim6\xi_c$ and failed when $L$ was much larger, where $\xi_c$ is the correlation depth scale of the infinitely wide random network with the same variances. Close to the critical line, where $\xi_c$ diverges, the trainable depth grew accordingly.
\end{observation}

How robust is this? The divergence of the trainable depth at the critical line has been reproduced many times; \textcite{XiaoEtAl2018} trained vanilla convolutional networks with $10{,}000$ layers, without residual connections or normalization, by initializing at criticality with orthogonal convolution kernels. The constant $6$ is empirical and depends on the setting, the training budget, and the criterion for success.\sidenote{The theory says nothing about the constant: it predicts a depth scale, not a training outcome.} Trainable is also not the same as trains well: \textcite{PenningtonEtAl2017} found that critical ReLU networks learn more slowly than critical sigmoidal networks with orthogonal weights, because criticality fixes the mean squared singular value of the input--output Jacobian and not its spread. Finally, modern networks escape the constraint by architecture rather than by tuning: residual connections and normalization layers make depth easy for reasons that the plain model does not contain (\cref{sec:signalprop:wrong}).

The questions for the model are therefore: What computes $\xi_c$? Why is there a sharp boundary, and why does the depth scale diverge on it? Where does the ReLU network at He initialization sit?

\section{The model: infinitely wide random networks}\label{sec:signalprop:model}

\begin{assumption}[Random fully connected network]\label[assumption]{ass:signalprop:net}
The input is $x\in\R^d$, the widths are $m_0=d$ and $m_1,\dots,m_L$, and the activation $\phi\colon\R\to\R$ is continuous and polynomially bounded: $\abs{\phi(t)}\le C_\phi(1+\abs t)^{k_\phi}$ for constants $C_\phi,k_\phi$. The \term{pre-activations}\index{pre-activation} are
\[
  h^{1}(x)=W^{1}x+b^{1},\qquad h^{\ell+1}(x)=W^{\ell+1}\phi\bigl(h^\ell(x)\bigr)+b^{\ell+1}\quad(1\le\ell<L),
\]
with $\phi$ applied coordinatewise, and all weights and biases are independent, with $W^{\ell}_{ij}\sim\cN(0,\sigma_w^2/m_{\ell-1})$ and $b^\ell_i\sim\cN(0,\sigma_b^2)$, where $\sigma_w^2>0$ and $\sigma_b^2\ge0$.
\end{assumption}

We write $a^0(x)=x$ and $a^\ell(x)=\phi(h^\ell(x))$ for the post-activations, so that $h^{\ell+1}=W^{\ell+1}a^\ell+b^{\ell+1}$ for $0\le\ell<L$. The widths are the hidden widths $m$ of the global notation, one per layer; the limit of the model sends them to infinity at fixed depth.

\paragraph{The kernel recursion.} The infinitely wide network is described by one deterministic function per layer, its \term{kernel}\index{kernel!of a random network}:
\begin{equation}\label{eq:signalprop:kernel}
  K^1(x,x')=\sigma_b^2+\sigma_w^2\frac{\ip{x,x'}}d,\qquad K^{\ell+1}(x,x')=\sigma_b^2+\sigma_w^2\,\E\bigl[\phi(z)\phi(z')\bigr],
\end{equation}
where $(z,z')$ is a centered Gaussian pair with $\E z^2=K^\ell(x,x)$, $\E z'^2=K^\ell(x',x')$, and $\E zz'=K^\ell(x,x')$. The polynomial bound makes the expectation finite. The two \term{order parameters}\index{order parameters!of signal propagation} are the length $q^\ell(x)=K^\ell(x,x)$, the limit of $\frac1{m_\ell}\norm{h^\ell(x)}^2$, and the correlation
\[
  c^\ell(x,x')=\frac{K^\ell(x,x')}{\sqrt{q^\ell(x)\,q^\ell(x')}},
\]
the limit of the cosine of the angle between $h^\ell(x)$ and $h^\ell(x')$. With $z\sim\cN(0,1)$, and $(u,v)$ standard Gaussians with correlation $c$, the length obeys the one-dimensional \term{length map}\index{length map}
\begin{equation}\label{eq:signalprop:length}
  q^{\ell+1}=\signalpropV(q^\ell),\qquad \signalpropV(q)=\sigma_b^2+\sigma_w^2\,\E\bigl[\phi(\sqrt q\,z)^2\bigr],
\end{equation}
and if both inputs sit at a fixed point $\bar q>0$ of $\signalpropV$, then the correlation obeys the \term{correlation map}\index{correlation map}
\begin{equation}\label{eq:signalprop:cmap}
  c^{\ell+1}=\signalpropC(c^\ell),\qquad \signalpropC(c)=\frac{\sigma_b^2+\sigma_w^2\,\E\bigl[\phi(\sqrt{\bar q}\,u)\phi(\sqrt{\bar q}\,v)\bigr]}{\bar q}.
\end{equation}
The fixed point is written $\bar q$, not $q^\star$, because a superscript star marks an optimum in this book. The quantity that will decide everything is the slope
\begin{equation}\label{eq:signalprop:chi}
  \chi_1=\sigma_w^2\,\E\bigl[\phi'(\sqrt{\bar q}\,z)^2\bigr].
\end{equation}

\begin{dictionary}
width $m_\ell\to\infty$ at fixed depth & width much larger than depth, $L/m\ll1$\\
weight and bias variances $\sigma_w^2,\sigma_b^2$ & the initialization scheme (LeCun, Glorot, He)\\
length $q^\ell$ & mean squared pre-activation of layer $\ell$\\
correlation $c^\ell(x,x')$ & cosine similarity of the representations of two inputs\\
fixed point $\bar c$ of $\signalpropC$ & how similar all inputs look to a deep layer\\
slope $\chi_1$ & growth factor per layer of squared gradient norms\\
depth scale $\xi$ & depth beyond which training fails (about $6\xi$ for $\tanh$ in \cref{obs:signalprop:depth})\\
critical line $\chi_1=1$ & the family of initializations that train deep networks\\
kernel $K^L$ & prior covariance of a wide Bayesian network; the NNGP kernel\\
random weights & weights at initialization only; training moves them\\
\end{dictionary}

\paragraph{What was thrown away, and why first.} The network is never trained: trainability is decided at initialization, at least in the sense of \cref{obs:signalprop:depth}, because a signal that does not reach the output cannot be learned from. The width is infinite, which turns sums over neurons into expectations and makes the theory deterministic, as the field of a magnet becomes deterministic when each spin interacts with infinitely many others. The architecture is a plain stack of dense layers, because it is the case in which depth is hardest. The data enter only through the inner products $\ip{x,x'}/d$, so two inputs are as good as a dataset.

\paragraph{What one had to notice.} Three things. Conditionally on one layer, the next is exactly Gaussian, because it is linear in Gaussian weights; only its covariance is random, and that covariance is an average over the neurons of the previous layer. Two inputs suffice: their lengths and their correlation are closed under the layer map, like the magnetization of a mean-field magnet. And the correlation map is a power series with nonnegative coefficients, so its behavior on $[0,1]$ is decided by its slope at the fixed point $c=1$.

\section{The solution: Gaussian processes at infinite width}\label{sec:signalprop:gp}

Let $\cX$ be a finite set of inputs, and write $h^\ell_i(\cX)=(h^\ell_i(x))_{x\in\cX}\in\R^{\cX}$ for the values of neuron $i$ of layer $\ell$ on these inputs. Let $\cF_\ell$ be the $\sigma$-algebra generated by $W^1,b^1,\dots,W^\ell,b^\ell$, with $\cF_0$ trivial.

\begin{lemma}[Exact conditional Gaussianity]\label[lemma]{lem:signalprop:conditional}
Under \cref{ass:signalprop:net}, for $0\le\ell<L$, conditionally on $\cF_\ell$, the vectors $h^{\ell+1}_1(\cX),\dots,h^{\ell+1}_{m_{\ell+1}}(\cX)$ are independent with law $\cN(0,\signalpropG^\ell)$, where
\[
  \signalpropG^\ell(x,x')=\sigma_b^2+\frac{\sigma_w^2}{m_\ell}\sum_{j=1}^{m_\ell}a^\ell_j(x)\,a^\ell_j(x'),\qquad x,x'\in\cX.
\]
In particular $\signalpropG^0=K^1$ on $\cX$, and the vectors $h^1_i(\cX)$ are independent with law $\cN(0,K^1)$.
\end{lemma}

\begin{proof}
The post-activations $a^\ell(x)$ are $\cF_\ell$-measurable, and the row $(W^{\ell+1}_{i1},\dots,W^{\ell+1}_{im_\ell},b^{\ell+1}_i)$ is a Gaussian vector independent of $\cF_\ell$ and of the other rows. Conditionally on $\cF_\ell$, the vector $h^{\ell+1}_i(\cX)$ is therefore a fixed linear image of this row, hence centered Gaussian, and the vectors for different $i$ are independent. Since the entries of the row are independent with variances $\sigma_w^2/m_\ell$ and $\sigma_b^2$, the conditional covariance of $h^{\ell+1}_i(x)=\sum_jW^{\ell+1}_{ij}a^\ell_j(x)+b^{\ell+1}_i$ and $h^{\ell+1}_i(x')$ is $\frac{\sigma_w^2}{m_\ell}\sum_ja^\ell_j(x)a^\ell_j(x')+\sigma_b^2$. For $\ell=0$, $a^0(x)=x$ and $m_0=d$ give $K^1$.
\end{proof}

The lemma is exact at every width. It says that a layer acts on a finite set of inputs through the $\abs\cX\times\abs\cX$ matrix $\signalpropG^\ell$ alone: in law, $h^{\ell+1}(\cX)$ is an $m_{\ell+1}\times\abs\cX$ matrix of independent standard Gaussians times $(\signalpropG^\ell)^{1/2}$. On two inputs, a network of width $m$ is therefore a Markov chain on $m\times2$ matrices, which can be simulated in time $O(mL)$ instead of $O(m^2L)$; all the finite-width simulations of this chapter use this reduction.\sidenote{We checked the reduction against simulations with full $m\times m$ weight matrices at $m=200$; the averaged correlations agree to within the Monte Carlo error.}

\begin{theorem}[Gaussian-process limit, sequential widths]\label{thm:signalprop:gp}\index{Gaussian process limit}
Under \cref{ass:signalprop:net}, let $\cX$ be a finite set of inputs, $1\le\ell\le L$, and $J\ge1$, with $m_\ell\ge J$. In the iterated limit in which $m_1\to\infty$, then $m_2\to\infty$, and so on up to $m_{\ell-1}\to\infty$, the law of $(h^\ell_1(\cX),\dots,h^\ell_J(\cX))$ converges weakly to the law of $J$ independent $\cN(0,K^\ell)$ vectors, with $K^\ell$ restricted to $\cX$. That is, for every bounded continuous $F\colon(\R^{\cX})^J\to\R$,
\[
  \lim_{m_{\ell-1}\to\infty}\cdots\lim_{m_1\to\infty}\E F\bigl(h^\ell_1(\cX),\dots,h^\ell_J(\cX)\bigr)=\E F(z_1,\dots,z_J),\qquad z_1,\dots,z_J\overset{\text{i.i.d.}}{\sim}\cN(0,K^\ell).
\]
For $\ell=1$ no limit is needed, and the law of $(h^\ell_1(\cX),\dots,h^\ell_J(\cX))$ does not depend on $m_\ell,\dots,m_L$.
\end{theorem}

\begin{proof}
The plan is induction on $\ell$: by \cref{lem:signalprop:conditional}, layer $\ell+1$ is a Gaussian whose covariance is a continuous function of the neurons of layer $\ell$; the induction hypothesis makes those neurons i.i.d.\ Gaussian, and the law of large numbers makes the covariance deterministic.

\emph{Two facts.} First, centered Gaussian laws depend continuously on their covariance: if $\Sigma_n\to\Sigma$, then the characteristic functions $\exp(-\frac12t^\top\Sigma_nt)$ converge to $\exp(-\frac12t^\top\Sigma t)$, and L\'evy's continuity theorem gives $\cN(0,\Sigma_n)\to\cN(0,\Sigma)$ weakly. Second, the last claim of the theorem holds because $h^\ell$ is a function of $W^1,b^1,\dots,W^\ell,b^\ell$, and the first $J$ neurons of layer $\ell$ use only the first $J$ rows of $W^\ell$ and $b^\ell$.

\emph{Base case.} By \cref{lem:signalprop:conditional} with $\ell=0$, the vectors $h^1_i(\cX)$ are i.i.d.\ $\cN(0,K^1)$ at every width.

\emph{Induction step.} Assume the statement for layer $\ell$ and every number of neurons, and fix $J$, $F$, and $m_\ell$. For $y=(y_1,\dots,y_{m_\ell})\in(\R^{\cX})^{m_\ell}$ define the $\cX\times\cX$ matrix $G(y)=\sigma_b^2\mathbf 1\mathbf 1^\top+\frac{\sigma_w^2}{m_\ell}\sum_j\phi(y_j)\phi(y_j)^\top$, and let $\Phi_F(y)=\E F(g_1,\dots,g_J)$ with $g_1,\dots,g_J$ i.i.d.\ $\cN(0,G(y))$. The function $\Phi_F$ is bounded by $\sup\abs F$, and it is continuous: $\phi$ is continuous, so $G$ is continuous, and the first fact applies. By \cref{lem:signalprop:conditional}, conditionally on $\cF_\ell$ the first $J$ neurons of layer $\ell+1$ are i.i.d.\ $\cN(0,G(h^\ell_1(\cX),\dots,h^\ell_{m_\ell}(\cX)))$, so
\[
  \E F\bigl(h^{\ell+1}_1(\cX),\dots,h^{\ell+1}_J(\cX)\bigr)=\E\,\Phi_F\bigl(h^\ell_1(\cX),\dots,h^\ell_{m_\ell}(\cX)\bigr).
\]
The induction hypothesis with $m_\ell$ neurons in place of $J$, applied to the bounded continuous function $\Phi_F$, shows that after the limits $m_1\to\infty,\dots,m_{\ell-1}\to\infty$ the right side becomes $\E\,\Phi_F(z_1,\dots,z_{m_\ell})$ with $z_j$ i.i.d.\ $\cN(0,K^\ell)$ (for $\ell=1$ this is the base case, with no limit).

Now let $m_\ell\to\infty$. Each entry of $\phi(z_j)\phi(z_j)^\top$ is integrable, since $\abs{\phi(s)\phi(t)}\le\frac12(\phi(s)^2+\phi(t)^2)$ and $\phi(z_j(x))^2$ is bounded by a polynomial in a Gaussian variable. By the strong law of large numbers, $G(z_1,\dots,z_{m_\ell})\to\sigma_b^2\mathbf 1\mathbf 1^\top+\sigma_w^2\E[\phi(z)\phi(z)^\top]$ almost surely, and by \eqref{eq:signalprop:kernel} the limit is $K^{\ell+1}$ on $\cX$. By the first fact, $\Phi_F(z_1,\dots,z_{m_\ell})\to\E F(z'_1,\dots,z'_J)$ almost surely, with $z'_i$ i.i.d.\ $\cN(0,K^{\ell+1})$, and bounded convergence gives the same limit for the expectation. This is the statement for layer $\ell+1$.
\end{proof}

The proof shows more: in the same iterated limit, now including $m_\ell\to\infty$, the empirical kernel $\signalpropG^\ell$ converges in probability to $K^{\ell+1}$, and in particular $\frac1{m_\ell}\norm{h^\ell(x)}^2\to q^\ell(x)$ (\cref{ex:signalprop:empirical}). This is the precise sense of \emph{mean-field exactness}: each neuron of layer $\ell+1$ feels the other neurons of layer $\ell$ only through an average over infinitely many of them, and the average does not fluctuate.

Taking the widths to infinity one at a time is a convenience of the proof. The theorem holds when they grow together.

\begin{theorem}[Simultaneous limits; \textcite{MatthewsEtAl2018}, \textcite{Hanin2023}]\label{thm:signalprop:simultaneous}
The conclusion of \cref{thm:signalprop:gp} holds when $\min(m_1,\dots,m_{\ell-1})\to\infty$ in an arbitrary manner, for activations satisfying a growth condition: \textcite{MatthewsEtAl2018} assume a linear envelope $\abs{\phi(t)}\le c+M\abs t$, and \textcite{Hanin2023} allows general nonlinearities and non-Gaussian weights under moment conditions.
\end{theorem}

For one hidden layer the statement is the central limit theorem, and \textcite{Neal1996} used it to show that Bayesian networks with a Gaussian prior on many hidden units define Gaussian-process priors; \textcite{LeeEtAl2018} and \textcite{MatthewsEtAl2018} extended it to depth, and \textcite{Yang2019} to most modern architectures.

\paragraph{The length map.} The kernel recursion needs a fixed length to define the correlation map. For odd activations such as $\tanh$ the length map has a unique attracting fixed point, and the proof is a Fourier argument.

\begin{proposition}[Length map of $\tanh$]\label[proposition]{prop:signalprop:length}
Let $\phi=\tanh$. The length map $\signalpropV$ of \eqref{eq:signalprop:length} is increasing and strictly concave on $[0,\infty)$, with $\signalpropV(0)=\sigma_b^2$, $\signalpropV'(0)=\sigma_w^2$, and $\signalpropV<\sigma_b^2+\sigma_w^2$.
\begin{enumerate}
\item If $\sigma_b^2>0$, then $\signalpropV$ has a unique fixed point $\bar q>0$, and $q^\ell\to\bar q$ monotonically from every $q^1\ge0$.
\item If $\sigma_b^2=0$ and $\sigma_w^2\le1$, then $q^\ell\to\bar q=0$ from every $q^1\ge0$.
\item If $\sigma_b^2=0$ and $\sigma_w^2>1$, then $\signalpropV$ has a unique positive fixed point $\bar q$, and $q^\ell\to\bar q$ monotonically from every $q^1>0$.
\end{enumerate}
In cases (i) and (iii), $\signalpropV'(\bar q)<1$, so the convergence is geometric.
\end{proposition}

\begin{proof}
The plan is to write $1-\tanh^2=\operatorname{sech}^2$ as a mixture of cosines with positive weights, which turns $q\mapsto\E[\operatorname{sech}^2(\sqrt qz)]$ into a mixture of decaying exponentials.

\emph{The representation.} The function $\operatorname{sech}^2$ is integrable, with Fourier transform $\int_\R\operatorname{sech}^2(t)\cos(\omega t)\dd t=g(\omega)$, where $g(\omega)=\pi\omega/\sinh(\pi\omega/2)$ for $\omega\ne0$ and $g(0)=2$ (\cref{ex:signalprop:fourier}). The function $g$ is positive and decays exponentially, so Fourier inversion gives $\operatorname{sech}^2(t)=\frac1{2\pi}\int_\R g(\omega)\cos(\omega t)\dd\omega$. Since $\E\cos(\omega\sqrt q z)=e^{-q\omega^2/2}$, Fubini's theorem gives
\[
  f(q)\defeq\E\bigl[\tanh^2(\sqrt q\,z)\bigr]=1-\frac1{2\pi}\int_\R g(\omega)\,e^{-q\omega^2/2}\dd\omega .
\]
\emph{Monotonicity and concavity.} The functions $\omega^2g(\omega)$ and $\omega^4g(\omega)$ are integrable, so we may differentiate under the integral for $q\ge0$:
\[
  f'(q)=\frac1{4\pi}\int_\R\omega^2g(\omega)e^{-q\omega^2/2}\dd\omega>0,\qquad f''(q)=-\frac1{8\pi}\int_\R\omega^4g(\omega)e^{-q\omega^2/2}\dd\omega<0 .
\]
At $q=0$, differentiating the inversion formula twice at $t=0$ gives $\frac1{2\pi}\int\omega^2g=-(\operatorname{sech}^2)''(0)=2$, because $\operatorname{sech}^2t=1-t^2+O(t^4)$; so $f'(0)=1$. With $\signalpropV=\sigma_b^2+\sigma_w^2f$ and $0\le f<1$, this proves the first sentence.

\emph{Case (ii).} For $q>0$, strict concavity gives $\signalpropV(q)<\signalpropV(0)+\signalpropV'(0)q=\sigma_w^2q\le q$. So $q^\ell$ decreases, its limit is a fixed point of the continuous map $\signalpropV$, and $0$ is the only one.

\emph{Cases (i) and (iii).} Let $g_0(q)=\signalpropV(q)-q$, strictly concave, with $g_0(q)\to-\infty$ because $\signalpropV$ is bounded. In case (i), $g_0(0)=\sigma_b^2>0$; in case (iii), $g_0(0)=0$ and $g_0'(0)=\sigma_w^2-1>0$; in both cases $g_0$ is positive somewhere on $(0,\infty)$ and has a zero $\bar q>0$ there by the intermediate value theorem. If $0<q_1<q_2$ were two zeros, then strict concavity on $[0,q_2]$ would give $g_0(q_1)>(1-q_1/q_2)g_0(0)\ge0$, a contradiction; so $\bar q$ is unique, $g_0>0$ on $(0,\bar q)$ and $g_0<0$ on $(\bar q,\infty)$. For $0<q<\bar q$, monotonicity gives $q<\signalpropV(q)<\signalpropV(\bar q)=\bar q$, and for $q>\bar q$ it gives $\bar q<\signalpropV(q)<q$: the iterates move monotonically toward $\bar q$, and their limit, a positive fixed point, is $\bar q$. In case (i) the start $q^1=0$ is mapped to $\sigma_b^2>0$. Finally, strict concavity gives $g_0'(\bar q)<\bigl(g_0(\bar q)-g_0(0)\bigr)/\bar q\le0$, that is, $\signalpropV'(\bar q)<1$.
\end{proof}

In case (ii) the network forgets the scale of its input: every length decays to zero, slowly when $\sigma_w^2=1$, since $f(q)=q-2q^2+O(q^3)$ gives $q^\ell\sim1/(2\ell)$. ReLU is the opposite extreme: $\E[\relu(\sqrt qz)^2]=q/2$, so with $\sigma_b=0$ the length map is the multiplication by $\sigma_w^2/2$, and every length is a fixed point exactly when $\sigma_w^2=2$.

\section{The solution: order, chaos, and the critical line}\label{sec:signalprop:phase}

Fix a fixed point $\bar q>0$ of the length map and two inputs with $q^1(x)=q^1(x')=\bar q$, so that $q^\ell=\bar q$ in every layer. Write $\bar\phi(t)=\phi(\sqrt{\bar q}\,t)$ and expand it in the normalized Hermite polynomials, with coefficients $a_k=\E[\bar\phi(z)h_k(z)]$ (\cref{lem:toolkit:hermite}). By \cref{cor:toolkit:dual}, the correlation map \eqref{eq:signalprop:cmap} is the power series
\begin{equation}\label{eq:signalprop:series}
  \signalpropC(c)=\sum_{k\ge0}\beta_kc^k,\qquad \beta_0=\frac{\sigma_b^2+\sigma_w^2a_0^2}{\bar q},\quad \beta_k=\frac{\sigma_w^2a_k^2}{\bar q}\ (k\ge1),
\end{equation}
with nonnegative coefficients, convergent on $[-1,1]$, with $\sum_k\beta_k=\signalpropC(1)=\signalpropV(\bar q)/\bar q=1$.\sidenote{The dual-activation view of \textcite{DanielyEtAl2016} is exactly this: a layer acts on correlations through the power series of its activation's squared Hermite coefficients, and composition of layers is composition of power series.}

\begin{theorem}[Order and chaos]\label{thm:signalprop:phase}\index{order--chaos transition}
Under \cref{ass:signalprop:net}, let $\phi$ be absolutely continuous and not constant, with $\phi'$ polynomially bounded, let $\bar q>0$ be a fixed point of $\signalpropV$, and let $\chi_1$ be given by \eqref{eq:signalprop:chi}.
\begin{enumerate}
\item The map $\signalpropC$ is continuous on $[-1,1]$, maps $[-1,1]$ into itself and $[0,1]$ into $[0,1]$, is nondecreasing and convex on $[0,1]$, satisfies $\signalpropC(1)=1$, and has left derivative $\signalpropC'(1)=\chi_1$ at $1$. For every $c\in[-1,1]$,
\begin{equation}\label{eq:signalprop:contraction}
  0\le1-\signalpropC(c)\le\chi_1\,(1-c).
\end{equation}
\item \emph{Ordered phase, $\chi_1<1$.} For every $c^1\in[-1,1]$, $1-c^\ell\le2\chi_1^{\ell-1}$, so $c^\ell\to1$. If $c^1\in(-1,1)$, then $c^\ell<1$ for every $\ell$ and $\frac1\ell\log(1-c^\ell)\to\log\chi_1$.
\item \emph{Critical line, $\chi_1=1$.} If $\phi$ is not affine and $c^1\in[0,1)$, then $c^\ell$ increases to $1$, and $(1-c^{\ell+1})/(1-c^\ell)\to1$: the decay is slower than every geometric rate.
\item \emph{Chaotic phase, $\chi_1>1$.} The map $\signalpropC$ has a unique fixed point $\bar c\in[0,1)$, and $\signalpropC'(\bar c)<1$. For every $c^1\in[0,1)$, $c^\ell\to\bar c$ monotonically, and $\abs{c^\ell-\bar c}=O(\rho^\ell)$ for every $\rho\in(\signalpropC'(\bar c),1)$.
\end{enumerate}
\end{theorem}

\begin{proof}
The plan is to read every claim off the power series \eqref{eq:signalprop:series}, using the identity
\begin{equation}\label{eq:signalprop:telescope}
  1-\signalpropC(c)=\sum_{k\ge1}\beta_k(1-c^k)=(1-c)\sum_{k\ge1}\beta_k\bigl(1+c+\dots+c^{k-1}\bigr),
\end{equation}
which follows from $\sum_k\beta_k=1$.

\emph{Part (i).} The coefficients are nonnegative and summable, so the series converges uniformly on $[-1,1]$ and $\signalpropC$ is continuous there; on $[0,1]$ every term is nondecreasing and convex. Since $\abs{c^k}\le1$, the first expression in \eqref{eq:signalprop:telescope} is nonnegative, and $\signalpropC(c)\ge\beta_0-\sum_{k\ge1}\beta_k\ge-1$; for $c\in[0,1]$ also $\signalpropC(c)\ge\beta_0\ge0$. For the slope, the Hermite coefficients of $\bar\phi'$ are $\sqrt{k}\,a_k$ by Gaussian integration by parts (\cref{lem:toolkit:hermite}, as in the proof of \cref{lem:toolkit:price}), so Parseval's identity and $\bar\phi'(t)=\sqrt{\bar q}\,\phi'(\sqrt{\bar q}\,t)$ give
\[
  \sum_{k\ge1}k\beta_k=\frac{\sigma_w^2}{\bar q}\sum_{k\ge1}ka_k^2=\frac{\sigma_w^2}{\bar q}\E\bigl[\bar\phi'(z)^2\bigr]=\sigma_w^2\E\bigl[\phi'(\sqrt{\bar q}\,z)^2\bigr]=\chi_1 ,
\]
which is finite because $\phi'$ is polynomially bounded. For $c\in[0,1)$ the sum in \eqref{eq:signalprop:telescope} increases to $\sum_kk\beta_k$ as $c\uparrow1$, by monotone convergence, so the left derivative at $1$ is $\chi_1$. For $c\in[-1,1]$, $\abs{1+c+\dots+c^{k-1}}\le k$ bounds the sum by $\chi_1$, which proves \eqref{eq:signalprop:contraction}.

\emph{Part (ii).} Iterating \eqref{eq:signalprop:contraction} gives $1-c^\ell\le\chi_1^{\ell-1}(1-c^1)\le2\chi_1^{\ell-1}$. Since $\phi$ is not constant, $\beta_k>0$ for some $k\ge1$, so for $c\in(-1,1)$ the first expression in \eqref{eq:signalprop:telescope} is positive and $\signalpropC(c)\ge\beta_0-\sum_{k\ge1}\beta_k\abs c^k>-1$: the interval $(-1,1)$ is invariant, and $c^\ell<1$ for every $\ell$. The ratio $(1-c^{\ell+1})/(1-c^\ell)$ is the sum in \eqref{eq:signalprop:telescope} at $c=c^\ell$. As $c^\ell\to1$, each term $\beta_k(1+c^\ell+\dots+(c^\ell)^{k-1})$ tends to $k\beta_k$ and is bounded by $k\beta_k$, which is summable; by dominated convergence the ratio tends to $\chi_1$, and the logarithms of the ratios have Ces\`aro means converging to $\log\chi_1$, which is the claim.

\emph{Part (iii).} If $\phi$ is not affine, then $\beta_k>0$ for some $k\ge2$ (if $a_k=0$ for all $k\ge2$, then $\bar\phi$ is affine), so for $c\in[0,1)$ the sum in \eqref{eq:signalprop:telescope} is strictly less than $\sum_kk\beta_k=1$, and $\signalpropC(c)>c$. Since $\signalpropC$ maps $[0,1]$ into itself, $c^\ell$ increases to a fixed point in $(c^1,1]$, and $\signalpropC(c)>c$ on $[0,1)$ leaves only $1$. The ratio tends to $\chi_1=1$ by the argument of part (ii).

\emph{Part (iv).} Let $g_1(c)=\signalpropC(c)-c$, convex and continuous on $[0,1]$, with $g_1(0)=\beta_0\ge0$, $g_1(1)=0$, and left derivative $\chi_1-1>0$ at $1$; so $g_1<0$ on some interval $(1-\delta,1)$. By the intermediate value theorem $g_1$ has a zero in $[0,1-\delta]$; let $\bar c$ be the largest one. On $(\bar c,1)$, $g_1$ has no zero and takes negative values, so $g_1<0$ there. If $0\le c_0<\bar c$ were another zero, then convexity on $[c_0,1]$ with $g_1(c_0)=g_1(1)=0$ would give $g_1\le0$ on $[c_0,1]$, and the convex function $g_1$ would have the interior maximum $g_1(\bar c)=0$ on this interval, forcing $g_1\equiv0$ on $[c_0,1]$, which contradicts $g_1<0$ near $1$. So $\bar c$ is the unique fixed point in $[0,1)$, $g_1>0$ on $[0,\bar c)$, and $g_1<0$ on $(\bar c,1)$. Because $\signalpropC$ is nondecreasing on $[0,1]$, $c<\bar c$ implies $c<\signalpropC(c)\le\bar c$, and $c>\bar c$ implies $\bar c\le\signalpropC(c)<c$; so $c^\ell$ moves monotonically toward $\bar c$ and converges to a fixed point in $[0,1)$, which is $\bar c$. By convexity, $\signalpropC'(\bar c)\le\bigl(\signalpropC(c)-\signalpropC(\bar c)\bigr)/(c-\bar c)<1$ for each $c\in(\bar c,1)$, where the last inequality is $g_1(c)<0$. Finally, $\signalpropC$ is differentiable at $\bar c<1$, so for $\rho\in(\signalpropC'(\bar c),1)$ there is $\eta>0$ with $\abs{\signalpropC(c)-\bar c}\le\rho\abs{c-\bar c}$ for $\abs{c-\bar c}\le\eta$; once $c^\ell$ is within $\eta$ of $\bar c$ the distance decreases by the factor $\rho$ per layer.
\end{proof}

\Cref{thm:signalprop:phase} says that a deep random network has two fates. In the ordered phase all inputs end up with the same representation at a geometric rate, and the network computes, at depth, a function that ignores its input. In the chaotic phase two inputs, however close, end up with correlation $\bar c<1$, a correlation that the network assigns to every pair of distinct inputs: nearby points are pulled apart, and at depth the network behaves like a random function of its input. By Price's theorem (\cref{lem:toolkit:price}), $\signalpropC'(c)=\sigma_w^2\E[\phi'(\sqrt{\bar q}u)\phi'(\sqrt{\bar q}v)]$, so the slope at $\bar c$ is a correlation of derivatives, as $\chi_1$ is a mean square of derivatives.

The rate of approach to the stable fixed point defines the \term{depth scale}\index{depth scale}
\begin{equation}\label{eq:signalprop:xi}
  \xi=-\frac1{\log\signalpropC'(\bar c)},\qquad \bar c=1\ \text{in the ordered phase, so that}\ \xi=-\frac1{\log\chi_1}\ \text{there}.
\end{equation}
It is the $\xi_c$ of \cref{obs:signalprop:depth}. The \term{critical line}\index{critical line} is the curve $\chi_1=1$ in the plane $(\sigma_w^2,\sigma_b^2)$; for $\tanh$ it starts at $(1,0)$, where $\bar q=0$, and $\chi_1$ is evaluated at the fixed length $\bar q(\sigma_w^2,\sigma_b^2)$ of \cref{prop:signalprop:length}. How does $\xi$ diverge as the line is approached?

\begin{heuristic}[Critical exponent of the depth scale]
Let $\signalpropC''(1)=\sum_kk(k-1)\beta_k=\sigma_w^2\bar q\,\E[\phi''(\sqrt{\bar q}\,z)^2]$ be finite and positive, and expand near $c=1$: $\signalpropC(1-\eps)=1-\chi_1\eps+\frac12\signalpropC''(1)\eps^2+O(\eps^3)$. In the ordered phase $\xi=-1/\log\chi_1\approx1/(1-\chi_1)$ as $\chi_1\uparrow1$. In the chaotic phase, the fixed point $\bar c=1-\bar\eps$ solves $\chi_1\bar\eps\approx\bar\eps+\frac12\signalpropC''(1)\bar\eps^2$, so $\bar\eps\approx2(\chi_1-1)/\signalpropC''(1)$, and $\signalpropC'(\bar c)\approx\chi_1-\signalpropC''(1)\bar\eps\approx2-\chi_1$. Both sides give
\[
  \xi\approx\frac1{\abs{\chi_1-1}}\qquad(\chi_1\to1),
\]
a divergence with exponent $1$, the value of mean-field theory.\qedhere
\end{heuristic}

The derivation assumes that $\bar c\to1$ as $\chi_1\downarrow1$, which fails at the end point $\sigma_b=0$ of the $\tanh$ line: there $\signalpropC$ is odd, $\bar c=0$ throughout the chaotic phase, and the transition is degenerate (\cref{ex:signalprop:exponent}). The numerics of \cref{fig:signalprop:phase}(b) confirm the formula on both sides at $\sigma_b^2=0.05$: at $\chi_1=0.998$, $\xi=496$ against $1/(1-\chi_1)=497$, and at $\chi_1=1.0029$, $\xi=349$ against $345$.

\paragraph{ReLU at He initialization.} For $\phi=\relu$ and $\sigma_b=0$, the length map is $q\mapsto\frac{\sigma_w^2}2q$, so lengths neither explode nor vanish exactly when $\sigma_w^2=2$, the initialization of \textcite{HeEtAl2015}. Then $\chi_1=2\,\E[\ind\{z>0\}]=1$: He initialization is critical. The correlation map of the arc-cosine kernel of \textcite{ChoSaul2009} decides how a critical ReLU network forgets.

\begin{theorem}[ReLU at He initialization]\label{thm:signalprop:relu}\index{ReLU network!He initialization}
Let $\phi=\relu$, $\sigma_w^2=2$, and $\sigma_b=0$. Then $\signalpropV(q)=q$ for every $q\ge0$, and for every $\bar q>0$ the correlation map is
\begin{equation}\label{eq:signalprop:arccos}
  \signalpropC(c)=\frac{\sqrt{1-c^2}+(\pi-\arccos c)\,c}{\pi},\qquad c\in[-1,1],
\end{equation}
with $\signalpropC'(c)=1-\frac{\arccos c}\pi$. For $\eps\in[0,1]$,
\begin{equation}\label{eq:signalprop:relu-expansion}
  \signalpropC(1-\eps)=1-\eps+\frac{2\sqrt2}{3\pi}\eps^{3/2}+R(\eps),\qquad 0\le R(\eps)\le\frac{\pi-2}{5\sqrt2\,\pi}\eps^{5/2}.
\end{equation}
For every $c^1\in[-1,1)$, $c^\ell$ increases to $1$ from $\ell=2$ on, and
\begin{equation}\label{eq:signalprop:ninepi}
  \lim_{\ell\to\infty}\ell^2\,(1-c^\ell)=\frac{9\pi^2}2 .
\end{equation}
\end{theorem}

\begin{proof}
The plan is to compute $\signalpropC$ by integrating its derivative from $c=0$, then to expand near $c=1$, and finally to linearize the recursion in the variable $(1-c)^{-1/2}$.

\emph{Lengths.} By symmetry of $z$, $\E[\relu(\sqrt qz)^2]=\frac q2\E z^2=\frac q2$, so $\signalpropV(q)=2\cdot\frac q2=q$.

\emph{The map.} By homogeneity $\E[\relu(\sqrt{\bar q}u)\relu(\sqrt{\bar q}v)]=\bar q\,\kappa(c)$ with $\kappa(c)=\E[\relu(u)\relu(v)]$, so $\signalpropC=2\kappa$. Price's theorem (\cref{lem:toolkit:price}) applies to $\relu$, which is absolutely continuous with derivative $\ind\{t>0\}$, and gives $\kappa'(c)=\Pr[u>0,v>0]$ for $c\in(-1,1)$. Write $u=\ip{e,g}$ and $v=\ip{e',g}$ for a standard Gaussian $g\in\R^2$ and unit vectors $e,e'$ at angle $\theta=\arccos c$. The direction of $g$ is uniform on the circle, and the set of directions with both inner products positive is an arc of length $\pi-\theta$, so $\kappa'(c)=(\pi-\arccos c)/(2\pi)$. At $c=0$, independence gives $\kappa(0)=(\E\relu(u))^2=1/(2\pi)$. Integrating with the antiderivative $s\arccos s-\sqrt{1-s^2}$ of $\arccos s$,
\[
  \kappa(c)=\frac1{2\pi}+\frac1{2\pi}\int_0^c(\pi-\arccos s)\dd s=\frac{\sqrt{1-c^2}+(\pi-\arccos c)\,c}{2\pi},
\]
first for $c\in(-1,1)$ and then on $[-1,1]$ by continuity (\cref{cor:toolkit:dual}). This is \eqref{eq:signalprop:arccos}, and $\signalpropC'=2\kappa'$.

\emph{The expansion.} For $s\in[0,1]$ the series $\arcsin s=s+\sum_{n\ge1}\gamma_ns^{2n+1}$ has nonnegative coefficients, so $(\arcsin s-s)/s^3$ is nondecreasing and $0\le\arcsin s-s\le(\frac\pi2-1)s^3$. Since $\arccos(1-t)=2\arcsin\sqrt{t/2}$ for $t\in[0,2]$,
\[
  \arccos(1-t)=\sqrt{2t}+\varrho(t),\qquad 0\le\varrho(t)\le\frac{\pi-2}{2\sqrt2}\,t^{3/2}.
\]
Integrating $\signalpropC'$ from $1-\eps$ to $1$ and substituting $s=1-t$,
\[
  1-\signalpropC(1-\eps)=\int_0^\eps\Bigl(1-\frac{\arccos(1-t)}\pi\Bigr)\dd t=\eps-\frac{2\sqrt2}{3\pi}\eps^{3/2}-\frac1\pi\int_0^\eps\varrho(t)\dd t ,
\]
which is \eqref{eq:signalprop:relu-expansion} with $R(\eps)=\frac1\pi\int_0^\eps\varrho$, and the bound on $\varrho$ integrates to the bound on $R$.

\emph{Convergence.} For $c<1$, $\signalpropC'<1$ on $(c,1)$, so $1-\signalpropC(c)<1-c$, that is, $\signalpropC(c)>c$; also $\signalpropC\le1$, and $\signalpropC(c)\ge\signalpropC(-1)=0$ because $\signalpropC$ is nondecreasing. So $c^2\in[0,1)$, the sequence $(c^\ell)_{\ell\ge2}$ increases, and its limit is a fixed point of $\signalpropC$, of which $1$ is the only one.

\emph{The rate.} Let $\eps_\ell=1-c^\ell\to0$, $\eps_\ell>0$, and $u_\ell=\eps_\ell^{-1/2}$; write $a=\frac{2\sqrt2}{3\pi}$. By \eqref{eq:signalprop:relu-expansion}, $\eps_{\ell+1}=\eps_\ell(1-\delta_\ell)$ with $\delta_\ell=a\eps_\ell^{1/2}+R(\eps_\ell)/\eps_\ell$, and $0\le\delta_\ell-a\eps_\ell^{1/2}\le\eps_\ell^{3/2}$. For $\delta\in[0,\frac12]$, Taylor's theorem gives $(1-\delta)^{-1/2}=1+\frac\delta2+r(\delta)$ with $0\le r(\delta)\le2\delta^2$. Once $\delta_\ell\le\frac12$,
\[
  u_{\ell+1}-u_\ell=u_\ell\Bigl(\frac{\delta_\ell}2+r(\delta_\ell)\Bigr)=\frac a2+O\bigl(\eps_\ell^{1/2}\bigr),
\]
because $u_\ell\delta_\ell=a+O(\eps_\ell)$ and $u_\ell\delta_\ell^2=O(\eps_\ell^{1/2})$. Hence $u_{\ell+1}-u_\ell\to a/2$, and by the Stolz--Ces\`aro theorem $u_\ell/\ell\to a/2$. So $\ell^2\eps_\ell=(\ell/u_\ell)^2\to4/a^2=\frac{4\cdot9\pi^2}{8}=\frac{9\pi^2}2$.
\end{proof}

ReLU networks at He initialization decorrelate polynomially: no depth scale, but a power law. The critical $\tanh$ network also decays as a power, but a different one: with $\signalpropC''(1)$ finite, $\eps_{\ell+1}=\eps_\ell-\frac12\signalpropC''(1)\eps_\ell^2+O(\eps_\ell^3)$ gives $1-c^\ell\sim2/(\signalpropC''(1)\ell)$ (\cref{ex:signalprop:smooth}). The exponent at the critical point depends on the smoothness of $\phi$ at the origin, through the first non-analytic term of $\signalpropC$ at $1$: $\eps^2$ for $\tanh$, $\eps^{3/2}$ for ReLU.

\paragraph{Gradients.} Forward correlations are half of the story; the other half is the size of the gradient. Let $\Loss$ be a loss that depends on the network through $h^L$, and let $\delta^\ell=\partial\Loss/\partial h^\ell$. The chain rule gives $\delta^\ell=D^\ell(W^{\ell+1})^\top\delta^{\ell+1}$ with $D^\ell=\diag(\phi'(h^\ell))$.

\begin{heuristic}[Gradient norms]
Assume the \term{gradient independence assumption}\index{gradient independence assumption}: in second moments, the matrix $W^{\ell+1}$ of the backward pass may be replaced by an independent copy, and $\delta^{\ell+1}$ is independent of $h^\ell$. Then, conditionally on $\delta^{\ell+1}$ and $h^\ell$, each coordinate $\sum_iW^{\ell+1}_{ij}\delta^{\ell+1}_i$ has variance $\frac{\sigma_w^2}{m_\ell}\norm{\delta^{\ell+1}}^2$, and summing over $j$,
\[
  \E\norm{\delta^\ell}^2=\sigma_w^2\,\E\Bigl[\frac1{m_\ell}\sum_{j}\phi'(h^\ell_j)^2\Bigr]\,\E\norm{\delta^{\ell+1}}^2\ \longrightarrow\ \chi_1\,\E\norm{\delta^{\ell+1}}^2
\]
at infinite width and length $\bar q$, by the law of large numbers of \cref{thm:signalprop:gp}. Iterating, $\E\norm{\delta^\ell}^2\approx\chi_1^{L-\ell}\,\E\norm{\delta^L}^2$: gradients vanish geometrically in the ordered phase, explode in the chaotic phase, and keep their mean squared size on the critical line, with depth scale $\xi_\nabla=-1/\log\chi_1$.\qedhere
\end{heuristic}

The assumption is false at finite width, where the forward and backward passes share the same matrices, but it gives the correct answer in the infinite-width limit for this network and many others: \textcite{Yang2020} gave a simple checkable condition under which computations with it are exact in the limit, and showed that they can be wrong when the condition fails.\sidenote{For $\tanh$ networks of depth $30$ at $\sigma_b^2=0.05$ we measured the ratio $\norm{\delta^0}^2/\norm{\delta^{30}}^2$. At criticality, averaged over $40$ networks, it was $0.89$ at width $800$ and $1.00$ at width $1600$, against $\chi_1^{30}=1$, with a spread across networks of about $25\%$. Over $10$ networks of width $800$ it was $0.0071$ against $0.0064$ at $\sigma_w^2=1.2$, and $265$ against $297$ at $\sigma_w^2=3$.} Two cases can be done exactly at every width, and they show what the mean hides.

\begin{theorem}[Jacobian norms at finite width]\label{thm:signalprop:finite}\index{finite-width fluctuations}
Let all widths equal $m$, $\sigma_b=0$, and fix an input $x\ne0$ and a direction $\dot x\in\R^d$. Let $u^\ell=\partial_th^\ell(x+t\dot x)|_{t=0}$, so that $u^1=W^1\dot x$ and $u^{\ell+1}=W^{\ell+1}D^\ell u^\ell$, and let $Y_\ell=\norm{D^\ell u^\ell}^2$, the squared norm of the derivative of $a^\ell=\phi(h^\ell)$ along $\dot x$.
\begin{enumerate}
\item If $\phi(t)=t$ (a deep linear network), then $Y_L=Y_1\prod_{\ell=1}^{L-1}F_\ell$, where $F_1,\dots,F_{L-1}$ are i.i.d., independent of $Y_1$, with law $\frac{\sigma_w^2}m\chi^2_m$.
\item If $\phi=\relu$ with $D^\ell=\diag(\ind\{h^\ell>0\})$ and $\sigma_w^2=2$, then the same holds with $F_\ell$ of law $\frac2m\chi^2_{B}$, where $B\sim\mathrm{Bin}(m,\frac12)$ and $\chi^2_0=0$.
\end{enumerate}
In both cases $\E Y_L=\chi_1^{L-1}\E Y_1$, and
\[
  \frac{\E Y_L^2}{(\E Y_L)^2}=\frac{\E Y_1^2}{(\E Y_1)^2}\Bigl(1+\frac\kappa m\Bigr)^{L-1},\qquad \kappa=2\ \text{(linear)},\quad \kappa=5\ \text{(ReLU)}.
\]
\end{theorem}

\begin{proof}
The plan is to show that, conditionally on $\cF_\ell$, the ratio $Y_{\ell+1}/Y_\ell$ has a fixed law that does not depend on $\cF_\ell$; independence of the factors and the moment formulas follow.

\emph{Linear case.} Here $D^\ell=I_m$ and $u^{\ell+1}=W^{\ell+1}u^\ell$. Conditionally on $\cF_\ell$, the coordinates of $W^{\ell+1}u^\ell$ are independent $\cN(0,\frac{\sigma_w^2}m\norm{u^\ell}^2)$, so $Y_{\ell+1}=Y_\ell\,F_\ell$ with $F_\ell\sim\frac{\sigma_w^2}m\chi^2_m$ independent of $\cF_\ell$ (if $Y_\ell=0$, then $Y_{\ell+1}=0$ and $F_\ell$ is drawn independently). Since $Y_1,F_1,\dots,F_{\ell-1}$ are $\cF_\ell$-measurable, the factors are independent of each other and of $Y_1$.

\emph{ReLU case.} Conditionally on $\cF_\ell$, \cref{lem:signalprop:conditional} (with $\cX=\{x\}$, extended to the pair $(h,u)$ by the same argument) shows that the pairs $(h^{\ell+1}_i,u^{\ell+1}_i)$, $i\le m$, are i.i.d.\ centered Gaussian, with $\Var(u^{\ell+1}_i)=\frac2m Y_\ell$ and $\Var(h^{\ell+1}_i)=\frac2m\norm{a^\ell}^2$. If $a^\ell=0$, then $D^\ell=0$, so $Y_\ell=0$ and $Y_{\ell+1}=0$. Otherwise $\Var(h^{\ell+1}_i)>0$, and the law of a centered Gaussian pair is invariant under $(h,u)\mapsto(-h,-u)$, so for every Borel set $A$, $\Pr[h>0,u^2\in A]=\Pr[h<0,u^2\in A]=\frac12\Pr[u^2\in A]$: the sign of $h^{\ell+1}_i$ is a fair coin independent of $(u^{\ell+1}_i)^2$. Hence $Y_{\ell+1}=\sum_i\ind\{h^{\ell+1}_i>0\}(u^{\ell+1}_i)^2$ is, conditionally on $\cF_\ell$, a sum of $B\sim\mathrm{Bin}(m,\frac12)$ independent squares of $\cN(0,\frac2mY_\ell)$ variables, that is, $Y_\ell$ times a variable of law $\frac2m\chi^2_B$. The case $\ell=1$ uses $x\ne0$, and independence follows as in the linear case.

\emph{Moments.} For the linear case $\E F=\sigma_w^2$ and $\E F^2=\frac{\sigma_w^4}{m^2}(m^2+2m)=\sigma_w^4(1+\frac2m)$, and $\chi_1=\sigma_w^2$. For ReLU, $\E[\chi^2_B\mid B]=B$ and $\E[(\chi^2_B)^2\mid B]=B^2+2B$, with $\E B=\frac m2$ and $\E B^2=\frac m4+\frac{m^2}4$; so $\E F=1=\chi_1$ and $\E F^2=\frac4{m^2}\bigl(\frac{m^2}4+\frac m4+m\bigr)=1+\frac5m$. Independence gives $\E Y_L=\E Y_1(\E F)^{L-1}$ and $\E Y_L^2=\E Y_1^2(\E F^2)^{L-1}$.
\end{proof}

At criticality the mean of $Y_L/Y_1$ is exactly $1$ at every width, as the heuristic predicts, but $\log(Y_L/Y_1)$ is a sum of $L-1$ independent terms, each with variance close to $\kappa/m$: the typical network has $\log(Y_L/Y_1)\approx-\frac{\kappa}{2m}L\pm\sqrt{\kappa L/m}$. The mean is carried by rare networks with large gradients, and the parameter that controls the typical behavior is $L/m$, which the infinite-width limit sets to zero. For ReLU this is the result of \textcite{Hanin2018} and \textcite{HaninNica2020}, who prove the analogous statements for the input--output Jacobian with variable widths, where $\sum_\ell1/m_\ell$ replaces $L/m$.

\paragraph{Dynamical isometry.} Criticality controls the mean squared singular value of the input--output Jacobian; trainability also depends on its spread. A network has \term{dynamical isometry}\index{dynamical isometry} if all singular values of its Jacobian concentrate near $1$. \Textcite{SaxeEtAl2014} observed that deep linear networks with orthogonal weight matrices have this property exactly, while Gaussian weights give a spread that grows with depth. \Textcite{PenningtonEtAl2017} computed the limiting singular value distribution with free probability: ReLU networks cannot achieve dynamical isometry with Gaussian or orthogonal weights, whereas $\tanh$ networks can, with orthogonal weights on the critical line near $\bar q=0$, and then learn faster.

\section{Reading the solution}\label{sec:signalprop:reading}

\paragraph{The phase diagram.} \Cref{fig:signalprop:phase}(a) shows the critical line of $\tanh$ networks, computed by solving $\chi_1=1$ with $\bar q$ from \cref{prop:signalprop:length}. It starts at $(\sigma_w^2,\sigma_b^2)=(1,0)$ and bends to larger weight variances as the bias variance grows: a bias adds a common component to all pre-activations, which raises $\bar q$, pushes the neurons into the saturated part of $\tanh$, and lowers $\E[\phi'^2]$, so more weight variance is needed to keep $\chi_1=1$. Left of the line is order and right of it is chaos, and the contours of $\xi$ crowd toward the line from both sides. Panel (b) shows $\xi$ along the cut $\sigma_b^2=0.05$, where the critical point is $\sigma_w^2=1.761$: the divergence is the critical slowing down of a second-order transition, and the curve $1/\abs{\chi_1-1}$ of the heuristic derivation is indistinguishable from it near the line. Combined with \cref{obs:signalprop:depth}, the plot is a map of trainable depth: at $\sigma_w^2=1.2$, where $\xi=5.9$, networks much deeper than $6\xi\approx36$ layers should fail; at $\sigma_w^2=1.7$, where $\xi=75$, several hundred layers should train.

\begin{figure}[t]
\begin{fullwidth}
\centering
\begin{tikzpicture}
\begin{groupplot}[group style={group size=2 by 1,horizontal sep=17mm},hplot,
  width=0.5\linewidth,height=0.37\linewidth,unbounded coords=jump]
\nextgroupplot[xlabel={$\sigma_w^2$},ylabel={$\sigma_b^2$},xmin=0,xmax=8,ymin=0,ymax=2,
  title={(a) phase diagram of $\tanh$},legend pos=north west,legend style={font=\scriptsize}]
\addplot[theory] table[x=crit,y=sb2] {figures/data/05-signalprop-phase.dat};
\addlegendentry{critical line $\chi_1=1$}
\addplot[theory2] table[x=ord20,y=sb2] {figures/data/05-signalprop-phase.dat};
\addlegendentry{$\xi=20$}
\addplot[theory3] table[x=ord5,y=sb2] {figures/data/05-signalprop-phase.dat};
\addlegendentry{$\xi=5$}
\addplot[theory2,forget plot] table[x=cha20,y=sb2] {figures/data/05-signalprop-phase.dat};
\addplot[theory3,forget plot] table[x=cha5,y=sb2] {figures/data/05-signalprop-phase.dat};
\node[font=\small,text=muted] at (axis cs:1.0,1.2) {ordered};
\node[font=\small,text=muted] at (axis cs:6.2,0.5) {chaotic};
\nextgroupplot[xlabel={$\sigma_w^2$},ylabel={depth scale $\xi$},xmin=0.5,xmax=8,ymode=log,ymin=1,ymax=2000,
  title={(b) $\xi$ along $\sigma_b^2=0.05$},legend pos=north east,legend style={font=\scriptsize},
  log ticks with fixed point]
\addplot[guide,forget plot] coordinates {(1.761,1) (1.761,2000)};
\addplot[theory] table[x=sw2,y=xi] {figures/data/05-signalprop-xi.dat};
\addlegendentry{$-1/\log\signalpropC'(\bar c)$}
\addplot[theory3,densely dashed] table[x=sw2,y expr={1/abs(\thisrow{chi1}-1)}] {figures/data/05-signalprop-xi.dat};
\addlegendentry{$1/\abs{\chi_1-1}$}
\end{groupplot}
\end{tikzpicture}
\caption{Infinitely wide $\tanh$ networks, computed by quadrature (theory only). (a)~The critical line $\chi_1=1$ in the plane of weight and bias variances, and the contours $\xi=5$ and $\xi=20$ of the depth scale \eqref{eq:signalprop:xi} on both sides of it. (b)~The depth scale along $\sigma_b^2=0.05$ (critical point $\sigma_w^2=1.761$, dotted), with the approximation $1/\abs{\chi_1-1}$ of the heuristic derivation.}
\label{fig:signalprop:phase}
\end{fullwidth}
\end{figure}

\paragraph{Correlations, layer by layer.} \Cref{fig:signalprop:corr} follows two inputs with initial correlation $0.6$ through $\tanh$ networks of width $1000$ in the three phases, against the iterated correlation map. In the ordered phase ($\sigma_w^2=1.2$, $\chi_1=0.845$) the correlation approaches $1$ geometrically, the gap $1-c^\ell$ shrinking by about $\chi_1$ per layer; at criticality it approaches $1$ like $6.2/\ell$, the power law of \cref{ex:signalprop:smooth}; in the chaotic phase ($\sigma_w^2=4$) it falls to $\bar c=0.166$. Every average lies within three standard errors of its curve. The spread across networks is the finite-width correction: it is small in the ordered phase, where the map contracts strongly, and largest in the chaotic phase, where the representation is a random function of the input.

\begin{figure}[t]
\centering
\begin{tikzpicture}
\begin{axis}[hplot,height=0.5\linewidth,xlabel={layer $\ell$},ylabel={correlation $c^\ell$},xmin=0,xmax=51,ymin=0,ymax=1.02,
  legend pos=south east,legend columns=1,legend style={font=\scriptsize}]
\addplot[guide,forget plot,domain=0:51,samples=2] {0.165578};
\addplot[theory] table[x=layer,y=ordth] {figures/data/05-signalprop-corr.dat};
\addlegendentry{ordered, $\sigma_w^2=1.2$}
\addplot[theory2] table[x=layer,y=critth] {figures/data/05-signalprop-corr.dat};
\addlegendentry{critical, $\sigma_w^2=1.761$}
\addplot[theory3] table[x=layer,y=chath] {figures/data/05-signalprop-corr.dat};
\addlegendentry{chaotic, $\sigma_w^2=4$}
\addplot[simulation,forget plot,each nth point=2] table[x=layer,y=ordmc] {figures/data/05-signalprop-corr.dat};
\addplot[simulation2,forget plot,each nth point=2] table[x=layer,y=critmc] {figures/data/05-signalprop-corr.dat};
\addplot[simulation,mark=triangle*,mark size=2pt,forget plot,each nth point=2] table[x=layer,y=chamc] {figures/data/05-signalprop-corr.dat};
\end{axis}
\end{tikzpicture}
\caption{Correlation of the pre-activations of two inputs in $\tanh$ networks with $\sigma_b^2=0.05$, starting from $c^1=0.6$ at the fixed length $\bar q$. Curves: iterates of the correlation map \eqref{eq:signalprop:cmap}. Marks: networks of width $m=1000$, averages over $50$ networks of the empirical cosine $\ip{h^\ell(x),h^\ell(x')}/(\norm{h^\ell(x)}\norm{h^\ell(x')})$, every second layer (seed $20261017$); the standard deviation across networks is about $0.03$ at criticality and $0.07$ in the chaotic phase. The dotted line is the chaotic fixed point $\bar c=0.166$.}
\label{fig:signalprop:corr}
\end{figure}

\paragraph{The ReLU constant.} \Cref{fig:signalprop:relu}(a) shows $\ell^2(1-c^\ell)$ for two orthogonal inputs of a ReLU network at He initialization. The limit $9\pi^2/2=44.41$ is approached slowly: the proof gives $u_{\ell+1}-u_\ell=\frac a2+\frac{3a^2}{8u_\ell}+O(u_\ell^{-2})$, hence $u_\ell=\frac a2\ell+\frac{3a}4\log\ell+O(1)$ and a relative correction of about $3\log\ell/\ell$ to $\ell^2(1-c^\ell)$; at $\ell=1000$ the exact iterate gives $43.2$, at $\ell=10^4$ it gives $44.26$. Simulations at width $1000$ agree with the iterated map up to depth $181$, within two standard errors.

\paragraph{Depth over width.} \Cref{fig:signalprop:relu}(b) tests \cref{thm:signalprop:finite} at width $m=50$ up to depth $151$: the variance of $\log(Y_L/Y_1)$ grows linearly in $L/m$, with slope close to $5$ for ReLU and $2$ for linear networks. At $L=m$ the typical gradient differs from its mean by a factor of order $e^{\pm\sqrt5}$, ten either way, although every network is critical on average. The same parameter controls the forward correlations: with the exact two-input chain at widths $25$ to $200$, the mean of $1-c^L$ at $L=m$ still matches $\signalpropC$, but its median is only about a quarter of the mean, at every width (\cref{ex:signalprop:fluct}).

\begin{figure}[t]
\begin{fullwidth}
\centering
\begin{tikzpicture}
\begin{groupplot}[group style={group size=2 by 1,horizontal sep=17mm},hplot,
  width=0.5\linewidth,height=0.37\linewidth,unbounded coords=jump]
\nextgroupplot[xlabel={layer $\ell$},ylabel={$\ell^2(1-c^\ell)$},xmode=log,xmin=1.5,xmax=1100,ymin=0,ymax=50,
  title={(a) ReLU, He initialization},legend pos=south east,legend style={font=\scriptsize},log ticks with fixed point]
\addplot[guide,domain=1.5:1100,samples=2] {44.4132};
\addlegendentry{$9\pi^2/2$}
\addplot[theory] table[x=layer,y=theory] {figures/data/05-signalprop-relu.dat};
\addlegendentry{iterated map}
\addplot[simulation,error bars/.cd,y dir=both,y explicit] table[x=layer,y=mc,y error=mcse] {figures/data/05-signalprop-relu.dat};
\addlegendentry{width $1000$}
\nextgroupplot[xlabel={depth over width $L/m$},ylabel={$\Var\log(Y_L/Y_1)$},xmin=0,xmax=3.1,ymin=0,ymax=17,
  title={(b) fluctuations at width $m=50$},legend pos=north west,legend style={font=\scriptsize}]
\addplot[theory] table[x=Lm,y=relutheory] {figures/data/05-signalprop-fluct.dat};
\addlegendentry{ReLU, theorem}
\addplot[theory2] table[x=Lm,y=lintheory] {figures/data/05-signalprop-fluct.dat};
\addlegendentry{linear, theorem}
\addplot[simulation] table[x=Lm,y=relumc] {figures/data/05-signalprop-fluct.dat};
\addlegendentry{ReLU, simulation}
\addplot[simulation2] table[x=Lm,y=linmc] {figures/data/05-signalprop-fluct.dat};
\addlegendentry{linear, simulation}
\end{groupplot}
\end{tikzpicture}
\caption{Critical networks at large and at finite width. (a)~ReLU at $\sigma_w^2=2$, $\sigma_b=0$, two orthogonal inputs: the iterated arc-cosine map \eqref{eq:signalprop:arccos} and its limit \eqref{eq:signalprop:ninepi}; points are averages over $400$ networks of width $1000$, with standard errors (seed $20261017$). (b)~Variance of the log squared Jacobian norm $\log(Y_L/Y_1)$ of \cref{thm:signalprop:finite}. Lines: $(L-1)\Var\log F$, computed exactly from the laws of the factors ($\Var\log F=0.108$ for ReLU and $0.041$ for linear networks at $m=50$). Points: $2000$ networks of width $50$ with full Gaussian weight matrices.}
\label{fig:signalprop:relu}
\end{fullwidth}
\end{figure}

\paragraph{The physics.} The order parameters $q$ and $c$ play the role of the magnetization of a mean-field magnet: at infinite width each neuron couples to infinitely many others, the coupling is through an average, and the average obeys a closed deterministic equation. Depth plays the role of time, or of the distance from a boundary at which a perturbation is applied, and $\xi$ is a correlation length. The ordered phase is the high-temperature phase, in which perturbations of the input die out, and the chaotic phase is the analogue of a positive Lyapunov exponent: infinitesimally close inputs separate by a factor $\sqrt{\chi_1}$ per layer, because $1-c^{\ell+1}\approx\chi_1(1-c^\ell)$ is the squared distance of the two representations in units of $\bar q$. At the critical line the correlation length diverges with the mean-field exponent $1$, and at the critical point the approach to equilibrium is a power law, as in critical slowing down. Mean-field theory is exact here for the same reason as in the Curie--Weiss model: the interaction range, the width, is infinite.

\section{What the model gets wrong}\label{sec:signalprop:wrong}

\paragraph{What was thrown away.} \emph{Finite width.} The limit sets $L/m=0$, and every fluctuation of order $\sqrt{L/m}$ disappears with it (\cref{thm:signalprop:finite}). \emph{Training.} The theory describes the network at time zero; it predicts where gradient descent can start, not what it finds, how fast, or whether the weights move far from their initial values (\cref{ch:lazy}). \emph{Architecture.} No residual connections, no normalization, no convolutions, no attention; each changes the recursion, and some remove the phase transition altogether. \emph{Data.} Only pairwise inner products of inputs enter, so the theory cannot distinguish a dataset from a set of points with the same Gram matrix, and it says nothing about which functions of the data are easy to learn. \emph{Weights.} Gaussian weights are not essential for the forward recursion \parencite{Hanin2023}, but orthogonal weights change the Jacobian spectrum and with it the trainability \parencite{PenningtonEtAl2017}.

\paragraph{Where it breaks.} \emph{Depth comparable to width.} Gradient statistics at finite width are controlled by $\sum_\ell1/m_\ell$, which is $L/m$ for constant widths: for ReLU networks the variance of the squared entries of the input--output Jacobian grows exponentially in this sum \parencite{Hanin2018}, and the logarithm of the squared norm of a Jacobian--vector product is approximately Gaussian, with mean and variance of order $\sum_\ell1/m_\ell$, when depth and width grow together \parencite{HaninNica2020}. \Cref{thm:signalprop:finite} is the exactly solvable shadow of this result. When $L\gtrsim m$ the infinite-width kernel is the mean of a broad distribution and not the typical value: the network that is actually initialized is not the Gaussian process. \Textcite{RobertsEtAl2022} organize the finite-width corrections as an expansion in $L/m$, and find that criticality is necessary to keep the corrections from growing exponentially, but that even at criticality the expansion parameter is $L/m$.

\emph{Criticality is necessary, not sufficient.} Critical ReLU networks train more slowly than critical sigmoidal networks with orthogonal weights, because the spread of the Jacobian singular values, which the model does not compute, grows with depth for ReLU \parencite{PenningtonEtAl2017}. Inside the trainable region, which initialization trains fastest depends on the activation beyond $\chi_1$ \parencite{HayouEtAl2019}. And training can leave the region in which the initialization analysis applies.

\emph{Kernels are not networks.} The kernel $K^L$ is the covariance of the infinitely wide Bayesian network \parencite{LeeEtAl2018}. In the large study of \textcite{LeeEtAl2020}, these kernels (and the tangent kernels of \cref{ch:lazy}) outperformed finite fully connected networks trained in the standard way, but underperformed finite convolutional networks; large learning rates and weight decay, which help finite networks, break the correspondence. The infinite-width kernel misses what trained convolutional networks gain by learning features.

\emph{Normalization changes the answer.} The model suggests that a well-chosen activation and variance cure exploding gradients. \Textcite{YangEtAl2019} showed that in fully connected networks with batch normalization and no skip connections, gradients explode exponentially with depth at initialization whatever the activation and the weight variance: the normalization itself causes the explosion, and such networks are untrainable at large depth.

\emph{Attention collapses.} A stack of self-attention layers without skip connections and feedforward blocks converges to a rank-one output doubly exponentially in depth \parencite{DongEtAl2021}, an ordered phase that variance tuning cannot leave. In transformers at initialization, this rank collapse makes the gradients of queries and keys vanish, and \textcite{NociEtAl2022} prevent it by scaling the residual branches with depth.

\paragraph{Fine structure.} The $1/m$ corrections are computed systematically by the effective theory of \textcite{RobertsEtAl2022}, in which the four-point vertex of the pre-activations grows linearly in $L/m$ at criticality. When depth and width go to infinity together with $L/m$ fixed, the covariance of a network whose activation is shaped toward the identity at a rate tied to the width converges to the solution of a stochastic differential equation, the \emph{neural covariance SDE} of \textcite{LiEtAl2022}: the deterministic recursion of this chapter becomes a diffusion, and $L/m$ is its time. Residual networks do not have the exponential depth scales of plain networks, and their signal propagation is subexponential or polynomial \parencite{YangSchoenholz2017}; scaling the residual branches by $1/\sqrt L$ gives well-defined infinite-depth limits in which hyperparameters transfer across depth \parencite{YangEtAl2024,BordelonEtAl2024}. Deep kernel shaping \parencite{MartensEtAl2021} controls the whole correlation map, not only its slope at $1$, by transforming the activation, and trains deep networks without skip connections or normalization. Open: a theory of trainability, rather than of signal propagation, that predicts the constant in $L\lesssim6\xi$, and a description of the regime $L\sim m$ during training.

\section{Model card}\label{sec:signalprop:card}

\modelcard{Deep networks train only for well-scaled initializations; the maximal trainable depth is a multiple of a depth scale that diverges on an order--chaos boundary; He initialization for ReLU.}
{Fully connected network with Gaussian weights $\cN(0,\sigma_w^2/m)$ and biases $\cN(0,\sigma_b^2)$, widths to infinity at fixed depth; order parameters $q^\ell$ and $c^\ell$.}
{Gaussian-process limit (\cref{thm:signalprop:gp}); length map (\cref{prop:signalprop:length}); order and chaos from the dual power series (\cref{thm:signalprop:phase}); ReLU at He initialization (\cref{thm:signalprop:relu}); exact Jacobian fluctuations at finite width (\cref{thm:signalprop:finite}).}
{The critical line $\chi_1=1$; depth scales $\xi=-1/\log\signalpropC'(\bar c)\approx1/\abs{\chi_1-1}$; gradient growth $\chi_1^{L}$; $1-c^\ell\sim9\pi^2/(2\ell^2)$ for ReLU and $\sim2/(\signalpropC''(1)\ell)$ for smooth critical activations.}
{Finite width (fluctuations of order $\sqrt{L/m}$); training; residual connections, normalization, and attention; the spread of the Jacobian spectrum; the constant in $L\lesssim6\xi$.}
{$1/m$ effective field theory; neural covariance SDE at fixed $L/m$; dynamical isometry; residual and depth-$\mu$P scalings; kernel shaping.}

The kernel $K^L$ of this chapter is the starting point of \cref{ch:lazy}: the neural tangent kernel is built from the same recursion together with the derivative map $c\mapsto\signalpropC'(c)$, and the order--chaos analysis decides whether that kernel is degenerate at depth. The Hermite expansion of the activation, which turns a layer into a power series in the correlation, returns in \cref{ch:singleindex}, where the lowest nonzero Hermite coefficient of a link function sets the information exponent; in both chapters Gaussian inputs make a nonlinearity act diagonally on Hermite polynomials, as characters diagonalize convolution in \cref{ch:grokking}. The spectrum of a product of random matrices in \cref{thm:signalprop:finite} is a relative of the Marchenko--Pastur law of \cref{ch:doubledescent}.

\begin{exercises}
\item \consolidation (The empirical kernel concentrates.) In the setting of \cref{thm:signalprop:gp}, show that in the iterated limit $m_1\to\infty,\dots,m_\ell\to\infty$ the empirical kernel $\signalpropG^\ell(x,x')$ converges in probability to $K^{\ell+1}(x,x')$, and deduce that $\frac1{m_\ell}\norm{h^\ell(x)}^2\to q^\ell(x)$. \hint{apply the theorem to bounded continuous functions of $\frac1J\sum_{j\le J}\phi(h^\ell_j(x))\phi(h^\ell_j(x'))$, truncated, and use the exchangeability of the neurons.}\label{ex:signalprop:empirical}

\item \consolidation (The Fourier transform of $\operatorname{sech}^2$.) (a) Integrate $e^{i\omega t}\operatorname{sech}^2t$, for $\omega>0$, over the boundary of the rectangle with corners $\pm R$ and $\pm R+i\pi$, and show that $\int_\R\operatorname{sech}^2(t)e^{i\omega t}\dd t=\pi\omega/\sinh(\pi\omega/2)$. (b) Deduce that $\E[\tanh^2(\sqrt qz)]=q-2q^2+O(q^3)$, and that for $\sigma_b=0$ and $\sigma_w^2=1$ the length of a $\tanh$ network satisfies $\ell q^\ell\to\frac12$. \hint{in (a), the only pole inside is the double pole at $i\pi/2$, where $\operatorname{sech}^2t\approx-(t-i\pi/2)^{-2}$, and the top side of the rectangle is $-e^{-\pi\omega}$ times the bottom one.}\label{ex:signalprop:fourier}

\item \consolidation\hard (The critical exponent, rigorously.) Assume $\signalpropC''(1)<\infty$. (a) Show that $\xi(1-\chi_1)\to1$ as $\chi_1\uparrow1$. (b) For $\tanh$ with a fixed $\sigma_b^2>0$, let $\sigma_w^2$ decrease to its critical value, and show that $1-\bar c\sim2(\chi_1-1)/\signalpropC''(1)$ and $\xi(\chi_1-1)\to1$ as $\chi_1\downarrow1$. (c) For $\tanh$ with $\sigma_b=0$, show that $\bar c=0$ for every $\sigma_w^2>1$, and that $\chi_1-1$ and $1-\signalpropC'(0)$ are both of order $(\sigma_w^2-1)^2$. Why does (b) not apply?\label{ex:signalprop:exponent}

\item \consolidation (Smooth activations at criticality.) Let $\chi_1=1$ and $0<\signalpropC''(1)<\infty$, and let $\signalpropC(1-\eps)=1-\eps+\frac12\signalpropC''(1)\eps^2+O(\eps^3)$. Show that $\ell(1-c^\ell)\to2/\signalpropC''(1)$, and that $\signalpropC''(1)=\sigma_w^2\bar q\,\E[\phi''(\sqrt{\bar q}z)^2]$. For $\tanh$ at $\sigma_b^2=0.05$, check numerically that $\signalpropC''(1)=0.324$, so that $1-c^\ell\approx6.17/\ell$, as in \cref{fig:signalprop:corr}. \hint{use $1/\eps_{\ell+1}-1/\eps_\ell$.}\label{ex:signalprop:smooth}

\item \consolidation\hard (The logarithmic correction.) In the proof of \cref{thm:signalprop:relu}, include the next term of the expansion \eqref{eq:signalprop:relu-expansion} and show that $u_\ell=\frac a2\ell+\frac{3a}4\log\ell+O(1)$. Deduce that $\ell^2(1-c^\ell)=\frac{9\pi^2}2\bigl(1-\frac{3\log\ell}\ell+O(\frac1\ell)\bigr)$, and compare with the values $43.2$ at $\ell=1000$ and $44.26$ at $\ell=10^4$.\label{ex:signalprop:log}

\item \variation (The error function.) Let $\phi=\erf$. (a) Show that for a centered Gaussian pair with variances $q_1,q_2$ and covariance $q_{12}$, $\E[\erf(z)\erf(z')]=\frac2\pi\arcsin\frac{2q_{12}}{\sqrt{(1+2q_1)(1+2q_2)}}$, the kernel of \textcite{Williams1997}. (b) Show that $\chi_1=\frac{4\sigma_w^2}{\pi\sqrt{1+4\bar q}}$, that the critical point at $\sigma_b=0$ is $\sigma_w^2=\pi/4$, and that the critical line is the curve $\sigma_w^2=\frac\pi4\sqrt{1+4\bar q}$, $\sigma_b^2=\bar q-\frac12\sqrt{1+4\bar q}\arcsin\frac{2\bar q}{1+2\bar q}$, $\bar q>0$. \hint{in (a), use Price's theorem and $\E[\phi'(z)\phi'(z')]$ for Gaussian $\phi'$.}\label{ex:signalprop:erf}

\item \variation (Leaky ReLU.) Let $\phi(t)=\max(t,\alpha t)$ with $\alpha\in[0,1)$ and $\sigma_b=0$. Show that lengths are preserved exactly when $\sigma_w^2=2/(1+\alpha^2)$, that this initialization is critical, and that
\[
  \signalpropC(c)=\frac{(1+\alpha)^2c+(1-\alpha)^2\frac2\pi\bigl(\sqrt{1-c^2}+c\arcsin c\bigr)}{2(1+\alpha^2)} .
\]
Deduce that $\ell^2(1-c^\ell)\to\frac{9\pi^2}2\bigl(\frac{1+\alpha^2}{(1-\alpha)^2}\bigr)^2$. Why does the constant blow up as $\alpha\to1$?\label{ex:signalprop:leaky}

\item \variation (ReLU off the He point.) Let $\phi=\relu$. (a) For $\sigma_w^2<2$ and $\sigma_b^2>0$, show that $\bar q=\sigma_b^2/(1-\sigma_w^2/2)$ and $\chi_1=\sigma_w^2/2<1$: ReLU networks with a fixed length are never chaotic. (b) For $\sigma_w^2>2$, show that lengths grow geometrically, so that \cref{thm:signalprop:phase} does not apply; compute the correlation map at $\sigma_b=0$ and show that it does not depend on $\sigma_w^2$.\label{ex:signalprop:relubias}

\item \variation (A residual network.) Let $h^{\ell+1}=h^\ell+\frac{\gamma}{\sqrt L}W^{\ell+1}\relu(h^\ell)$ with $W^{\ell+1}_{ij}\sim\cN(0,2/m)$ and no biases. (a) Derive the length recursion $q^{\ell+1}=(1+\gamma^2/L)q^\ell$ at infinite width, and show that $q^L/q^1\to e^{\gamma^2}$ as $L\to\infty$, while without the factor $1/\sqrt L$ the length grows like $(1+\gamma^2)^L$. (b) Derive the correlation recursion and show that, with the factor $1/\sqrt L$, the correlation of two inputs converges as $L\to\infty$ to the solution at time $1$ of an ordinary differential equation in $c$; write it down.\label{ex:signalprop:resnet}

\item \simulation (Gradients in three phases.) Build $\tanh$ networks of width $800$ and depth $30$ with $\sigma_b^2=0.05$ and $\sigma_w^2\in\{1.2,1.761,3\}$, start from a random input of length $\bar q$, and backpropagate $\Loss=\ip{v,h^{30}}$ with $v\sim\cN(0,I_m)$. Average $\norm{\delta^0}^2/\norm{\delta^{30}}^2$ over $10$ networks and compare with $\chi_1^{30}$ ($0.0064$, $1$, and $297$). Observe agreement within the spread across networks, which is $30$--$60\%$ of the mean. Repeat at width $200$ and observe the spread grow, as \cref{thm:signalprop:finite} suggests.\label{ex:signalprop:grad}

\item \simulation (Typical versus mean.) Use the two-input Markov chain of \cref{lem:signalprop:conditional} for ReLU at He initialization, with orthogonal inputs, widths $m\in\{25,50,100,200\}$, depths $L\in\{m/4,m/2,m\}$, and $1000$ networks each. (a) Observe that the mean of $1-\hat c^L$ agrees with the iterated map within a few percent in every case. (b) Observe that the ratio of the median of $1-\hat c^L$ to its mean is about $0.75$, $0.5$, and $0.23$ for $L/m=\frac14,\frac12,1$, independently of $m$, and that the relative standard deviation grows from about $0.8$ to about $2.2$. Explain why the controlling parameter is $L/m$.\label{ex:signalprop:fluct}

\item \research (Trainability beyond signal propagation.) \Cref{obs:signalprop:depth} relates the trainable depth to $6\xi$, with a constant fitted to experiments. Open: a solvable model of the first steps of training of a deep network near the critical line that predicts the trainable depth, including the constant and its dependence on the learning rate and the training budget.\label{ex:signalprop:researchtrain}

\item \research (The regime $L\sim m$.) The neural covariance SDE of \textcite{LiEtAl2022} describes shaped networks at initialization when $L/m$ is fixed. Open: an analogous limit for unshaped critical networks (ReLU at He initialization) that describes the law of the Jacobian spectrum and of $c^L$ at $L/m$ fixed, and its consequences for the first steps of gradient descent.\label{ex:signalprop:researchLm}
\end{exercises}

\notesandrefs
Wide random networks became Gaussian processes with \textcite{Neal1996}, who studied priors for Bayesian networks with one hidden layer, and \textcite{Williams1997} computed the covariance functions for error-function and Gaussian units. \Textcite{ChoSaul2009} introduced the arc-cosine kernels as the covariances of infinitely wide rectifier layers and composed them into multilayer kernels, and \textcite{DanielyEtAl2016} developed the duality between activations and their kernels, including the power-series form of the correlation map. The deep Gaussian-process limit was stated by \textcite{LeeEtAl2018}, who used it for Bayesian inference on MNIST and CIFAR-10, and proved by \textcite{MatthewsEtAl2018}; the sequential proof given here is the argument behind both, made exact by \cref{lem:signalprop:conditional}. \Textcite{Hanin2023} gives a short proof for general weight distributions, and \textcite{Yang2019} extends the limit to most modern architectures.

The order--chaos picture is due to \textcite{PooleEtAl2016}, who also showed that in the chaotic phase the length of a curve of inputs grows exponentially with depth, a measure of expressivity, and the depth scales and the link to trainability to \textcite{SchoenholzEtAl2017}, who also showed that dropout destroys the critical point. \Textcite{PenningtonEtAl2017} and \textcite{XiaoEtAl2018} added the spectrum of the Jacobian and orthogonal initializations, building on the orthogonal initializations of deep linear networks of \textcite{SaxeEtAl2014}. \Textcite{HayouEtAl2019} analyze the edge of chaos for general activations, including ReLU-like ones.

The exact product structure of \cref{thm:signalprop:finite} is elementary for linear networks; for ReLU it rests on the sign symmetry used by \textcite{Hanin2018}, whose analysis, and that of \textcite{HaninNica2020}, covers the full Jacobian and variable widths. The book of \textcite{RobertsEtAl2022} develops the whole subject, criticality included, as a perturbation theory in $1/m$.\sidenote{The Hermite form of the correlation map makes the order--chaos theorem a statement about power series with nonnegative coefficients summing to one, generating functions of probability distributions; \cref{thm:signalprop:phase}(ii)--(iv) is then the extinction and survival dichotomy of a branching process, with $\chi_1$ as the mean offspring number.}

\paragraph{Translating the literature.} \Textcite{PooleEtAl2016} and \textcite{SchoenholzEtAl2017} write $N_l$ for our width $m_\ell$, $q^l_{aa}$ and $q^l_{ab}$ for our $q^\ell$ and $K^\ell(x,x')$, $\mathcal V$ for the length map, $q^*$ and $c^*$ for our fixed points $\bar q$ and $\bar c$, $c^l_{ab}$ for our $c^\ell$, and $\mathcal Dz$ for the standard Gaussian measure; their $\chi_1$ is ours, their $\chi_c$ is our $\signalpropC'(\bar c)$, their $\xi_c=-1/\log\chi_c$ is our $\xi$, and their $\xi_\nabla$ is our gradient depth scale $-1/\log\chi_1$. \Textcite{LeeEtAl2018} write $K^l(x,x')$ as we do, with $n_l$ for the width. \Textcite{RobertsEtAl2022} write $n_\ell$ for the width, $C_W$ and $C_b$ for our $\sigma_w^2$ and $\sigma_b^2$, $K^{(\ell)}$ for the kernel, and use parallel and perpendicular susceptibilities $\chi_\parallel$ and $\chi_\perp$, which at a fixed point are our $\signalpropV'(\bar q)$ and our $\chi_1$; their criticality condition $\chi_\parallel=\chi_\perp=1$ is the critical line together with marginal stability of the length. \Textcite{HaninNica2020} write $n_j$ for the widths and $\beta=\sum_j1/n_j$, which is $L/m$ for constant widths.

\printchapterbib
\end{refsection}
